% Original fixed-tier evaluation fixture for reSearch Investigate Set P.
% Synthetic physics note written from scratch for evaluation (not a real paper).
\documentclass{article}
\usepackage{amsmath}

\title{Weiss Mean-Field Theory and Landau Expansion of the Ising Ferromagnet}
\author{S. Halvorsen}
\date{}

\begin{document}
\maketitle

\begin{abstract}
We apply Weiss mean-field theory to a nearest-neighbor Ising ferromagnet on a
$z$-coordinated lattice, derive the magnetization self-consistency equation,
construct the Landau free energy near $T_c$, and compute the zero-field
susceptibility on both sides of the transition.
\end{abstract}

\section{Spin Hamiltonian}\label{sec:model}

Consider $N$ classical Ising spins $s_i \in \{-1,+1\}$ on a regular lattice of
coordination number $z$ with ferromagnetic coupling $J > 0$ and uniform field $h$:
\begin{equation}\label{eq:hamiltonian}
  H = -J \sum_{\langle i,j \rangle} s_i s_j - h \sum_{i=1}^{N} s_i,
\end{equation}
where $\langle i,j \rangle$ runs once over each of the $N z / 2$ nearest-neighbor bonds.

\section{Mean-field self-consistency}\label{sec:self-consistency}

Writing $s_i = m + \delta s_i$ with uniform magnetization $m = \langle s_i \rangle$
and dropping terms quadratic in $\delta s_i \delta s_j$ reduces \eqref{eq:hamiltonian}
to the decoupled mean-field Hamiltonian
\[
  H_{\mathrm{MF}} = \frac{1}{2} N z J m^2 - (z J m + h)\sum_{i=1}^{N} s_i.
\]
At inverse temperature $\beta = (k_B T)^{-1}$, evaluating $\langle s_i \rangle_{\mathrm{MF}}$
yields the Weiss self-consistency equation
\begin{equation}\label{eq:mfeq}
  m = \tanh\bigl(\beta (z J m + h)\bigr).
\end{equation}
For $h = 0$, a nonzero solution emerges continuously when the slope at $m = 0$
reaches unity, locating the critical temperature at $k_B T_c = z J$.

\section{Landau free energy and susceptibility}\label{sec:landau}

Near $T_c$ and for weak $h$, we invert and expand Eq. (2) around $m = 0$ using
$\operatorname{artanh}(m) = m + m^3/3 + O(m^5)$. Following the standard variational
treatment of continuous transitions \cite[Section~3.1]{kerr2019ising}, the resulting
equation of state follows from minimizing the Landau free energy per spin,
\begin{equation}\label{eq:landau-f}
  f(m) = f_0 + \frac{1}{2} k_B (T - T_c) m^2 + \frac{1}{12} k_B T m^4 - h m,
\end{equation}
with $f_0 = -k_B T \ln 2$. Differentiating \eqref{eq:landau-f} gives
$k_B(T - T_c) m + \frac{1}{3} k_B T m^3 = h$; taking $\partial m / \partial h$ at
$h = 0$ with $m = 0$ ($T > T_c$) and $m^2 = 3(T_c - T)/T$ ($T < T_c$) gives
\begin{equation}\label{eq:chi}
  \chi(T) = \left.\frac{\partial m}{\partial h}\right|_{h=0} =
  \begin{cases}
    \bigl[k_B (T - T_c)\bigr]^{-1}, & T > T_c, \\
    \bigl[2 k_B (T_c - T)\bigr]^{-1}, & T < T_c.
  \end{cases}
\end{equation}
Both branches diverge as $|T - T_c|^{-1}$ with an amplitude ratio of $2$.

\appendix

\section{Cumulant expansion check}\label{sec:app-cumulant}

For a trial independent-spin Hamiltonian $H_0 = -\eta \sum_{i=1}^{N} s_i$, the
first-order cumulant inequality bounds the exact free energy $F$ from above:
\[
  F \leq F_0 + \langle H - H_0 \rangle_0
  = -N k_B T \ln\bigl[2\cosh(\beta \eta)\bigr]
    - \frac{1}{2} N z J m_0^2 + N(\eta - h) m_0,
\]
where $m_0 = \tanh(\beta \eta)$. Minimizing over $\eta$ enforces $\eta = z J m_0 + h$,
confirming that the mean-field free energy is a variational upper bound.

\section*{Acknowledgments}

This synthetic note was prepared solely as an evaluation fixture and uses no
external correspondence or funded research.

\begin{thebibliography}{9}
\bibitem{kerr2019ising} M. Kerr, \emph{Mean-Field Lattice Models}, Synthetic Lecture Notes (2019), doi:10.0000/eval.ising.2019.
\end{thebibliography}

\end{document}
